\section{Lecture 3: Quantum Error Correction and Stabilizers - 09 February 2026}
Outline:
\begin{enumerate}[label=(\Roman*)]
    \item Quantum Coding
    \item Operator Measurements to Syndromes
    \item Shor's 9-qubit code
    \item QEC Criteria
    \item Stabilizers
\end{enumerate}

Last time: \\
Operator sum representation:
\begin{equation}
    \mathcal{E}(\rho) = \sum_k E_k \rho E_k^\dagger \text{, where } \sum_k E_k^\dagger E_k = I
\end{equation}
This equation describes how quantum noise is modeled. 

\subsection{Quantum Coding}
\begin{definition}
    An $[[n,k]]$ quantum code $C$ is a $k$-qubit subspace of an $n$-qubit Hilbert space. 
\end{definition}
\begin{example}
    $n = 3, k=1$ defined by $\ket{0_L} = \ket{000}$ and $\ket{1_L} = \ket{111}$ \\
    Errors: Bit flip: $\mathcal{E}_{BF}(\rho) = (1-p) \rho + p X\rho X$ (on a single qubit) \\
    Transmit $a\ket{000} + b\ket{111}$, then 
    \begin{center}
    \begin{tabular}{ c | c | c }
    Receive & Probability & Decode\\
    \hline
    $a\ket{000} + b\ket{111}$ & $(1 -p)^3$ & $a\ket{0} + b\ket{1}$ \\  
    $a\ket{001} + b\ket{110}$ & $p(1-p)^2$ & $a\ket{0} + b\ket{1}$ \\  
    $a\ket{010} + b\ket{101}$ & $p(1-p)^2$ & $a\ket{0} + b\ket{1}$ \\  
    $a\ket{100} + b\ket{001}$ & $p(1-p)^2$ & $a\ket{0} + b\ket{1}$ \\  
    $a\ket{011} + b\ket{100}$ & $p^2(1-p)$ & $a\ket{1} + b\ket{0}$ \\  
    $a\ket{101} + b\ket{010}$ & $p^2(1-p)$ & $a\ket{1} + b\ket{0}$ \\  
    $a\ket{100} + b\ket{001}$ & $p^2(1-p)$  & $a\ket{1} + b\ket{0}$\\  
    $a\ket{111} + b\ket{000}$ & $p^3$ & $a\ket{1} + b\ket{0}$
    \end{tabular}
    \end{center}
    The last 4 cases are the errors. The probability in decoded state:
    $P_err= 3p^2(1-p) + p^3 = O(p^2)$. In the original case, we had $O(p)$. 
\end{example}



\subsection{Operator Measurement and Syndromes}
\begin{definition}
    Given $U$ (could act on multiple qubits not just 1) with eigenvalues $\pm 1$ and eigenstates $\ket{u_\pm}$, ``measuring $U$" is 
    \begin{center}
    \begin{quantikz}[row sep=0.6cm, column sep=0.8cm, slice all, slice titles=Slice \col,slice style=blue,slice label style ={inner sep=1pt,anchor=south west,rotate=40}]
    \lstick{$\ket{0}$} & \gate{H} & \ctrl{1} & \gate{H} & \meter{m} \\
    \lstick{$c_0 \ket{u_-} + c_1 \ket{u_+}$}         & \qwbundle{} & \gate{U} & \qwbundle{} & \rstick{$\ket{\phi}$} 
    \end{quantikz}
    \end{center}
    At Slice 1: 
    \begin{equation}
        \frac{1}{\sqrt{2}}(\ket{0} + \ket{1})  [c_0 \ket{u_+} + c_1\ket{u_-}]
    \end{equation}
    At Slice 2:  
    \begin{equation}
        \frac{\ket{0}}{\sqrt{2}}[c_0 \ket{u_+} + c_1\ket{u_-}] + \frac{\ket{1}}{\sqrt{2}} [c_0 \ket{u_+} - c_1\ket{u_-}]
    \end{equation}
    At Slice 3:
    \begin{equation}
        \frac{\ket{0} + \ket{1}}{2}[c_0 \ket{u_+} + c_1\ket{u_-}] + \frac{\ket{0} + \ket{1}}{\sqrt{2}} [c_0 \ket{u_+} - c_1\ket{u_-}] = \ket{0}c_0\ket{u_+} + \ket{1} c_1 \ket{u_-}
    \end{equation}
    Measure $m=0$, $\ket{\phi} = \ket{U_+}$; Measure $m=1$, $\ket{\phi} = \ket{U_-}$.\\
    We have segmented the Hilbert space into its two halves $\ket{u_-}$ and $\ket{u_+}$
\end{definition}
Error correction and Syndromes: 
\begin{enumerate}
    \item Measure ``Syndrome Ops"
    \item Apply a unitary ``Recover Ops" $R$
\end{enumerate}
Let $U_1 = Z_1Z_2=ZZI$ and $U_2 = Z_2Z_3 = IZZ$, then
\begin{center}
    \begin{tabular}{ c | c | c | c | c | c}
    Receive & Probability & Decode & $U_1=ZZI$ & $U_2=IZZ$ & Recovery $R$\\
    \hline
    $a\ket{000} + b\ket{111}$ & $(1 -p)^3$ & $a\ket{0} + b\ket{1}$ & 0 & 0 & $I$\\  
    $a\ket{001} + b\ket{110}$ & $p(1-p)^2$ & $a\ket{0} + b\ket{1}$ & 0 & 1 & $X_3$\\ 
    $a\ket{010} + b\ket{101}$ & $p(1-p)^2$ & $a\ket{0} + b\ket{1}$ & 1 & 1 &$X_2$\\
    $a\ket{100} + b\ket{001}$ & $p(1-p)^2$ & $a\ket{0} + b\ket{1}$ & 1 & 0 & $X_1$\\
    $a\ket{011} + b\ket{100}$ & $p^2(1-p)$ & $a\ket{1} + b\ket{0}$ & & & \\
    $a\ket{101} + b\ket{010}$ & $p^2(1-p)$ & $a\ket{1} + b\ket{0}$ & & & \\
    $a\ket{100} + b\ket{001}$ & $p^2(1-p)$  & $a\ket{1} + b\ket{0}$ & & & \\
    $a\ket{111} + b\ket{000}$ & $p^3$ & $a\ket{1} + b\ket{0}$ & & & \\
    \end{tabular}
\end{center}
\begin{definition}
    Fidelity: If the final result is a density matrix $\rho$, given the ideal state $\ket{\psi}$, the definition of fidelity is given by $F(\ket{\psi}, \rho) = \braket{\psi|\rho|\psi} \geq 1 - 3p^2(1-p)-p^3 = 1 - O(p^2)$
\end{definition}

Quantum Circuit: 
\begin{center}
\begin{quantikz}
 \wireoverride{n} &  \wireoverride{n} & \wireoverride{n} & \wireoverride{n} \lstick{$\ket{0}$} & \gate{H} & \ctrl{1} & \ctrl{2} & \gate{H} & \meter{m_0} & \cw & \cwbend{1} \\ 
\lstick{$a\ket{0} + b\ket{1}$} & \ctrl{1}  & \ctrl{2} & \gate{\mathcal{E}_{BF}} & \qw  & \gate{Z} & \qw  & \qw  & \qw & \qw & \gate[wires=3]{R} & \qw \\
\lstick{$\ket{0}$} & \targ{} &\qw & \gate{\mathcal{E}_{BF}} & \qw & \qw & \gate{Z} & \gate{Z} & \qw  &\qw  & & \qw \\
\lstick{$\ket{0}$} &   \qw & \targ{} & \gate{\mathcal{E}_{BF}} & \qw & \gate{Z} & \qw & \qw & \qw  & \qw & & \qw \\
 \wireoverride{n} &  \wireoverride{n} & \wireoverride{n} & \wireoverride{n} \lstick{$\ket{0}$} & \gate{H} & \ctrl{-1} & \qw & \ctrl{-2} & \gate{H} & \meter{m_1} & \cwbend{-1} \\ 
\end{quantikz}
\end{center}
The result will be identical to the input state encoded state right after the CNOTs. This is QEC for a bit-flip channel. 


\subsection{Shor's 9-qubit code}
In the previous section, we have an error-correcting code for bit-flips. We need a quantum circuit phase-flip error-correcting code. \\
Recall: $\mathcal{E}_{PF} = (1-p)\rho + p Z\rho Z$ \\
Recall: $HZH = X$, $HXH = Z$, $\ket{\pm} = \frac{1}{\sqrt{2}}(\ket{0} + \ket{1})$. \\
We can write 
\begin{equation}
    H \mathcal{E}_{PF}(H\rho H )H = \mathcal{E}_{BF}(\rho) 
\end{equation}
Define 
\begin{align}
    \ket{0_L} &= \ket{+++} \label{plus}\\ 
    \ket{1_L} &= \ket{---} \label{minus}
\end{align}
The errors are $\ket{+} \Leftrightarrow \ket{-}$. 
\begin{claim}
    Arbitrary single qubit errors are a combination of $= X,Z,XZ$ 
\end{claim}
Recall: $\mathcal{E}(\rho) = \sum_k E_k \rho E_k^\dagger$. \\
Let $\sigma_j = \{I,X,Y,Z\}$. Then it follows that any $2 \times 2$ matrix can be written as a linear combination of those 4 matrices. Therefore, $E_k = \sum_j c_{kj} \sigma_j$. Then we get 
\begin{equation}
    \mathcal{E}(\rho) = \sum_{k, j'} c_{kj}\sigma_j \rho \sigma_{j'}c_{kj'}^* = \sum_{jj'} \chi_{jj'}\sigma_j\rho \sigma_{j'}
\end{equation}
where $\chi_{jj'} = \sum_k c_{kj}c_{kj}^*$. This $\chi$-matrix representation is a density matrix.
\begin{example}
    Consider a small rotation: $R_x(2\epsilon) \simeq \ket{\psi} - i \epsilon X \ket{\psi}$ ($X$ is bit flip). We want to say either identity happened or a bit flip happen and nothing in between. This contradicts because this a real unitary. Why can we say that either identity or a bit flip when it clearly experiences a superposition? \\
    \begin{equation}
        \mathcal{E}(\rho) \approx (1- \epsilon^2)\rho - i \epsilon X\rho + i \epsilon \rho X + \epsilon^2 X\rho X
    \end{equation}
    $- i \epsilon X\rho + i \epsilon \rho X$ are the offending terms and are called ``coherent" error. We do not have to worry about them because \textit{they disappear in the syndrome measurement}. The syndrome measurement projects the code state into either having a identity or a bit flip. \\
    To emphasize: \textit{syndrome measurement project environment into definite error states!}
\end{example}

9 qubit code: $[[9,1]]$, encodes 1 logical qubit with 9 physical qubits \\
Code states:
\begin{align}
    \ket{0_L} &= \frac{(\ket{000} + \ket{111})^{\otimes 3}}{\sqrt{8}} \quad \text{``}+_L+_L+_L \text{"}\\ 
    \ket{1_L} &= \frac{(\ket{000} - \ket{111})^{\otimes 3}} {\sqrt{8}} \quad \text{``}-_L-_L-_L \text{"}
\end{align}
The extra layer of the ``+++" and ``---" gives \eqref{plus} and \eqref{minus}, First layer corrects bit flips and the second layer corrects phase flips. This is called a \textit{concatenated code} - we concatenated a phase flip code on top of a bit flip code. \\
This code corrects any single qubit error. It does so using syndrome operators (and recover operators given by the following: \\
Bit flips: 
\begin{align}
    &Z_1Z_2  \quad Z_2Z_3 \\
    &Z_4Z_5  \quad Z_5Z_6 \\
    &Z_7Z_8  \quad Z_8Z_9 \\
    \text{corresponds to }& ZZI \hspace*{0.45cm} IZZ \text{ from previous example}
\end{align}
Phase flips:
\begin{align}
    X_1X_2X_3X_4X_5X_6 \\
    X_4X_5X_6X_7X_8X_9 \\
\end{align}
\begin{example}
    Suppose we have a $Z_3$ error. This means our states are 
    \begin{align}
        \ket{0_L} &= (000 - 111)(000 + 111)(000 + 111) \\
        \ket{1_L} &= (000 + 111)(000 - 111)(000 - 111) 
    \end{align}
    Which one of the syndrome operators is gonna give us 1? \\
    $Z_1Z_2 : m = 0$, $Z_2Z_3 : m = 0$, similarly for rest of bit flips \\
    $X_1X_2X_3X_4X_5X_6: m = 1$, $X_4X_5X_6X_7X_8X_9: m = 1$ \\
    The phase flips are parity operator for signs. The bit flips are parity operator for bit values. 
\end{example}
(What $n$ and $\#$ of syndrome operators are minimal?) \\
\begin{fact}
    9-qubit code: Every single qubit error gives a unique syndrome bit string. \\
    Some two-qubit errors also can give a unique syndrome bit string
\end{fact}
Issues: 
\begin{itemize}
    \item What are efficient families of $[[n,k]]$ codes? We would like the rate $\frac{k}{n}$ as $n \rightarrow \infty$ to not go to zero, i.e. finite rate code families
    \item We also want to be able to compute on codes
    \item We want minimum possible complexity of encoding and decoding (example: locally testable codes - where we can figure out the syndrome, what error occurred, by only looking at a small window of qubits?)   
\end{itemize}



\subsection{Quantum Error Correction Criteria}
\begin{theorem}
    Let $C$ be a quantum code defined by the orthonormal states $\{\ket{\psi_k}\}$ and let $\mathcal{E}(\rho) = \sum_k E_k \rho E_k^\dagger$ \\
    Then $\exists$ quantum operation $R$ correcting error  $\epsilon$ on the the code space $C$ iff
    \begin{itemize}
        \item orthogonality: $\braket{\psi_l |E_j^\dagger E_k|\psi_l} = 0,  \forall j \neq k, \forall l$
        \item non-deformation: $\braket{\psi_l|E_k^\dagger E_k| \psi_l} = c_k, \forall l $
    \end{itemize}
\end{theorem}

\textit{Correction after class:} \\ 
If $\{\ket{\psi_k}\}$ is an orthonormal basis for the codespace (where $k$ represents logical states $\ket{0_L}, \ket{1_L}$, etc.), the Knill--Laflamme (KL) criteria for quantum error
correction codes is:
\begin{equation}
\braket{\psi_a | E_i^\dagger E_j | \psi_b} = \alpha_{ij}\delta_{ab}
\end{equation}
What This Tells Us (The ``Rules of the Game'') \\
This criteria reveals two specific requirements that must be met for every pair of basis states
$|\psi_a\rangle$ and $|\psi_b\rangle$:
\begin{enumerate}
\item \textbf{The ``orthogonality'' requirement ($a\neq b$):}

When $a\neq b$, the right side of the equation becomes zero:
\begin{equation}
\braket{\psi_a | E_i^\dagger E_j | \psi_b} = 0
\end{equation}
This means that an error $E_j$ acting on one basis state must result in a state that is
orthogonal to the state produced by error $E_i$ acting on a different basis state.
Essentially, the errors cannot ``cross-talk'' or smear the two distinct pieces of logical
information into each other.

\item \textbf{The ``non-deformation'' requirement ($a=b$):}

When $a=b$, the equation simplifies to:
\begin{equation}
\braket{\psi_a | E_i^\dagger E_j | \psi_a} = \alpha_{ij}
\end{equation}
Crucially, $\alpha_{ij}$ is a constant that does not depend on $a$. This implies that the
transition amplitude between two error states must be identical for every basis vector in the
code. If this value changed depending on whether you were in state $\ket{\psi_1}$ or
$\ket{\psi_2}$, the environment would effectively be ``measuring'' the state of the qubit.
By keeping $\alpha_{ij}$ independent of the basis, the error process remains ``blind'' to the
stored quantum information.
\end{enumerate}


\subsection{Stabilizers}
\begin{definition}
    The Pauli group $G_n$ on $n$ qubits is all $n$-fold tensor products of $\{X,Z,XZ\}$ and $\{\pm 1, \pm i\}$ 
\end{definition}
\begin{example}
    For $n=1$, $G = \{X, Y,Z, \ldots\}$
\end{example}
Note: $\forall g, h \in G_n$, either $gh=hg$ (commute) or $gh=-hg$ (anticommute)
\begin{definition}
    A stabilizer for  $\{\ket{\psi_l}\} \doteq V_s$, ($\ket{\psi_l}$ are like basis vectors, $V_s$ is the vector space) is the set \begin{equation}S = \{g \in G_n \,\, |\,\,  g\ket{\psi} = \ket{\psi}, \, \forall \ket{\psi} \in V_s\}\end{equation} By convention $-I \notin S$. 
\end{definition}
Note that $S$ is an Abelian group. 
\begin{example} Stabilizer examples:
    \begin{itemize}
        \item $V_s = \{\ket{00}\}$, then $S = \{II, ZZ, ZI, IZ\} = \langle IZ,ZI \rangle $ 
        \item $V_s = \{\frac{\ket{00} + \ket{11}}{\sqrt{2}}\}$, then $S = \langle XX,ZZ \rangle$
        \item $V_s = \{\ket{000}, \ket{111}\}$, then $S = \langle ZZI, IZZ \rangle$
    \end{itemize}
    Note that if $\dim(V_s) = 2^k$, then $k = n - (\text{min } \# \text{ gens for } S)$
    
\end{example}